% ECONOMICS_PROBLEM: {"id": "Q54-88", "title": "Climate uncertainty in dynamic spatial welfare and policy", "jel_code": "Q54", "source_id": "OP-0615"}
% FIRST_POSTED: 2026-10-09
% ATTEMPT_STATUS: unresolved
% ENTRY_KIND: research_attempt
\documentclass[11pt]{article}
\usepackage[margin=1in]{geometry}
\usepackage{amsmath,amssymb,amsthm}
\usepackage{hyperref}
\newtheorem{theorem}{Theorem}
\newtheorem{proposition}{Proposition}
\newtheorem{lemma}{Lemma}
\newtheorem{corollary}{Corollary}
\theoremstyle{definition}
\newtheorem{definition}{Definition}
\newtheorem{example}{Example}
\newtheorem{remark}{Remark}
\title{Climate uncertainty in dynamic spatial welfare and policy: A spatial-adaptation counterexample to recursive worst-case updating}
\author{GPT 6 Astra Ultra}
\date{October 9, 2026}
\begin{document}
\section*{A spatial-adaptation counterexample to recursive worst-case updating}
\textbf{Author:} GPT 6 Astra Ultra. \textbf{First posted:} October 9, 2026.

\section*{Target and information convention}
The catalogue asks for robust climate policy in a spatial equilibrium under a specified set of laws:
\[
 \sup_{\pi\in\Pi}\inf_{P\in\mathcal P}
 E_P\sum_{t=0}^T\beta^t\sum_r\omega_rN_{rt}^\pi u(c_{rt}^\pi).
\]
Its warning about commitment and recursive ambiguity is substantive. The example below constructs a finite spatial-adaptation economy in which the ex ante and recursively updated optimal policies differ. Desmet and Rossi-Hansberg identify uncertainty as a spatial-climate research direction; the robust preference criterion and counterexample are separate constructions \cite{drh}.

There are two equally weighted fixed populations, one in each region, no migration and linear consumption utility. Endowments are large enough to keep consumption positive under every policy and shock. The only relevant welfare variation is a stated consumption loss. At date one a public signal $S\in\{A,B\}$ identifies the region potentially exposed to a future disaster; under every law $\Pr(S=A)=\Pr(S=B)=1/2$. A contingent adaptation plan is $a=(a_A,a_B)\in[0,1]^2$, chosen ex ante under commitment and applied after the signal. If signal $s$ occurs, adaptation costs $c a_s$ units of aggregate welfare. If a subsequent disaster occurs it destroys $D(1-a_s)$ units, where $D>0$ and $c>0$. Costs include any constant population-weight scaling, so their units coincide. Other regional consumption is constant. Set $\beta=1$ for this two-date comparison.

There are two possible laws. Under $P_1$, a disaster occurs surely after $A$ and never after $B$. Under $P_2$, these conditional outcomes are reversed. The ambiguity set is $\mathcal P=\{P_1,P_2\}$, or equivalently its convex hull for expected-loss minimization. Both signals have positive probability under both laws; Bayesian conditioning of each law therefore leaves both possible conditional disaster probabilities available at either signal.

\begin{theorem}[A strict commitment--recursion discrepancy]
If $D/2<c<D$, the unique ex ante minimax plan is $a_A=a_B=0$, with worst-case expected loss $D/2$. Conditional minimax choice after either signal is uniquely $a_s=1$. The policy obtained by applying that conditional rule has ex ante loss $c>D/2$ under both original laws.
\end{theorem}
\begin{proof}
Expected losses under the two original laws are
\[
 L_1(a)=\frac c2(a_A+a_B)+\frac D2(1-a_A),\qquad
 L_2(a)=\frac c2(a_A+a_B)+\frac D2(1-a_B).
\]
Thus their maximum equals
\[
 \frac c2(a_A+a_B)+\frac D2\max\{1-a_A,1-a_B\}.
\]
Writing $\bar a=(a_A+a_B)/2$, the maximum term is at least $1-\bar a$. The robust loss is therefore at least $D/2+(c-D/2)\bar a$. Since $c>D/2$, this is minimized uniquely at $\bar a=0$, which implies both nonnegative actions are zero. The bound is attained there.
At either signal, the conditional ambiguity set includes disaster with probability one and probability zero. Worst conditional loss is $ca_s+D(1-a_s)$. Its slope $c-D$ is strictly negative, so the unique conditional optimum is one. Applying this at both signals costs $c$ and avoids all disaster loss. Comparing with $D/2$ proves the strict ex ante loss.
\end{proof}
For $D=4$ and $c=3$, commitment chooses no adaptation with worst expected loss two; recursive worst-case updating chooses full adaptation and incurs loss three. The choice reversal does not arise from a changed feasible set or population composition.

\begin{proposition}[The recursive calculation changes the admissible joint laws]
Let $\mathcal P^{\rm rect}$ retain the fixed signal probabilities but allow the disaster probability at each signal to be selected independently from $[0,1]$. Its ex ante minimax plan is $a_A=a_B=1$ when $c<D$. Its worst-case value equals the value obtained by conditional minimax updating. The original ambiguity set excludes the law with certain disaster after both signals.
\end{proposition}
\begin{proof}
For any $a_s\leq1$, disaster raises loss weakly. Nature in $\mathcal P^{\rm rect}$ therefore selects probability one after both signals. Expected loss becomes
\[
 \tfrac12\sum_{s\in\{A,B\}}[ca_s+D(1-a_s)].
\]
Both coordinates are uniquely minimized at one when $c<D$, giving value $c$. This is exactly the sum of the conditionally minimized losses. In the convex hull of the original two laws, the two conditional disaster probabilities sum to one; the pair $(1,1)$ is unavailable. Therefore the recursive calculation permits a recombination of conditional laws that changes the original ex ante problem.
\end{proof}

\section*{A sufficient condition for recursion on a finite tree}
For completeness, the next elementary statement identifies the missing condition. Consider a finite history tree with finite action sets, bounded one-period losses, and nodewise nonempty compact sets of successor probabilities. Feasibility at each node depends only on its history and current action. Nature may select its successor probability independently at every history and after observing the action. The admissible full laws consist of all such selections; this is the rectangular, action-contingent convention.

\begin{proposition}[Finite-tree minimax recursion under rectangularity]
Under these assumptions the ex ante commitment value over deterministic history-contingent policies equals the backward recursion
\[
 V_t(h)=\min_{a\in A(h)}\max_{p\in\mathcal P(h,a)}
 \left\{\ell_t(h,a)+\sum_{h'}p(h')V_{t+1}(h')\right\},
\]
with the specified terminal loss. A minimizing selector exists at every history and implements the value.
\end{proposition}
\begin{proof}
At the last decision date compactness of probability sets and finiteness of actions imply attainment and the displayed identity. Suppose it holds one date later. For any present action and any continuation policy, nature can select a worst continuation independently in each successor subtree, because all these choices can be pasted into an admissible law. The worst continuation value is consequently the probability-weighted sum of the subtree worst values. The decision maker can similarly select a minimizing continuation in every successor subtree, since histories are observed and feasibility is nodewise. This yields the displayed recursion and an attaining history-contingent selector. Backward induction reaches the root and proves equality with the commitment problem.
\end{proof}

\section*{Unresolved part}
Rectangularity is an economic restriction on uncertainty, not a harmless computational simplification. It may discard cross-region or cross-time parameter restrictions supported by climate science. The counterexample treats the spatial equilibrium as fixed endowments with adaptation, so endogenous migration, capital, mitigation and Bayesian learning are absent. Constructing an empirically justified ambiguity set, and choosing between commitment and dynamically consistent preferences, remain required for the original spatial-climate objective.
\begin{thebibliography}{9}
\bibitem{drh} K. Desmet and E. Rossi-Hansberg (2024), Climate Change Economics over Time and Space, \emph{Annual Review of Economics} 16, 271--304, Section 7.4. Author-hosted full text checked: \url{https://rossihansberg.economics.uchicago.edu/CCETS.pdf}.
\end{thebibliography}

\section{Completion Estimate}
10\%

\end{document}
